\section{Introduction}
\label{sec:intro}

The rampant spread of digital misinformation and synthetic media poses acute risks to information integrity worldwide. While automated fact verification has experienced rapid advancement in high-resource languages through benchmarks such as FEVER~\citep{thorne2018fever} and FEVEROUS~\citep{aly2021feverous}, low-resource and typologically diverse languages remain largely underserved. Among these, the Kazakh language---a Kipchak Turkic tongue spoken by over twenty million individuals across Central Asia---exhibits structural characteristics that expose the failure modes of standard Natural Language Processing (NLP) paradigms.

Automating fact verification for Kazakh reveals three fundamental bottlenecks:
\begin{enumerate}
    \item \textbf{Agglutinative Morphology and Lexical Sparsity}: Kazakh features an extensive, rule-governed agglutinative morphological system where nominal and verbal stems can accumulate long chains of inflectional and derivational affixes marking case, plural, possessive, mood, tense, voice, negation, and evidentiality. A single root such as \textit{жаз} (to write) produces hundreds of distinct surface forms, such as \textit{жазбағандықтан} (because [one] did not write) or \textit{жаздырылмаған} (which was not caused to be written). Standard subword tokenizers, including Byte-Pair Encoding (BPE) and WordPiece, are typically trained on multilingual corpora dominated by isolating or fusional languages; consequently, they split Kazakh morphemes into arbitrary, semantically fragmented tokens, destroying morphological boundaries and exacerbating out-of-vocabulary sparsity.
    \item \textbf{The Shortcut Trap in Fact Verification}: In conventional fact verification benchmarks, claims classified as \textsc{Not Enough Info} (NEI) often suffer from topical or lexical divergence from candidate evidence. Language models trained on these datasets quickly adopt superficial heuristics: if a claim exhibits high unigram overlap with the evidence corpus, the model predicts \textsc{Supported}; if lexical overlap is low, it predicts \textsc{NEI}~\citep{schuster2019fever}. In agglutinative languages, this shortcut failure is amplified: subtle changes in single derivational or polarity affixes can entirely invert veracity while retaining over 90\% of surface token overlap, inducing devastating failure rates when evaluated on real-world misinformation.
    \item \textbf{The Dual Threat of Synthetic Disinformation}: Contemporary information ecosystems are inundated with content co-authored or entirely hallucinated by Large Language Models (LLMs). Verifying textual statements in isolation from generative provenance leaves systems vulnerable to sophisticated, hallucinated synthetic prose that mimics encyclopedic tone. A complete epistemic defense requires simultaneous estimation of factual grounding and generative origin.
\end{enumerate}

To bridge this critical gap, we present \textbf{Kazakh-FEVER 3K}, the first human-annotated, morphologically-balanced fact verification benchmark for Kazakh. Kazakh-FEVER 3K comprises 3,000 claim-evidence pairs curated by native linguists, featuring an explicit \textsc{Hard NEI} protocol designed to disarm lexical overlap shortcuts. We further develop a hybrid evidence retrieval pipeline combining morpheme-aware BM25 with dense contrastive representations, and introduce a 2D Trust Matrix that unifies factual entailment with synthetic text detection.

Our primary contributions are summarized as follows:
\begin{itemize}
    \item \textbf{Kazakh-FEVER 3K Benchmark}: We curate and release the first human-validated fact verification benchmark for Kazakh, featuring 3,000 carefully verified claim-evidence pairs divided into canonical training, development, encyclopedic test, and emerging news evaluation splits.
    \item \textbf{Hard NEI Protocol}: We formulate a systematic curation methodology that preserves high surface overlap while altering modal, evidential, and temporal operators, compelling models to perform deep compositional and morphological reasoning rather than relying on lexical matching shortcuts.
    \item \textbf{Hybrid Retrieval and 2D Trust Matrix Framework}: We design a dual-stage pipeline combining morpheme-stemmed sparse BM25 and multilingual dense embeddings (\textsc{mContriever}) via Reciprocal Rank Fusion, paired with a morpheme-grounded NLI verifier and a continuous 2D Trust Matrix that provides calibrated assessment across factual truth and generative authenticity.
\end{itemize}
